% VI. Експериментални резултати и анализ.
% СКЕЛЕТ. Клетките се попълват СЛЕД измерването. Празна клетка е по-добра от
% правдоподобно число: правдоподобното число се цитира месеци по-късно като факт.
\section{Експериментални резултати и анализ}
\label{sec:results}

\subsection{Проведени измервания}

Измерванията са проведени по методиката от раздел~\ref{sec:method}. Отчетени са
изпълненията, удовлетворяващи условията за валидност, а отхвърлените са изброени заедно
с причината за отхвърлянето.
\TODO{брой проведени, брой валидни и брой отхвърлени изпълнения, с причините}

\subsection{Време за сходимост}

Времето за сходимост по топология и по стойност на интервала за отпадане е дадено в
табличен вид \tabref{tab:convergence}.

\begin{table}[H]
\caption{Време за сходимост след отказ на връзка}
\label{tab:convergence}
\centering
\fontsize{11}{13}\selectfont
\begin{tabular}{lccc}
\toprule
Топология & Диаметър & Интервал за отпадане & Време за сходимост \\
\midrule
\TODO{} & \TODO{} & \TODO{} & \TODO{} \\
\TODO{} & \TODO{} & \TODO{} & \TODO{} \\
\bottomrule
\end{tabular}
\end{table}

\TODO{коментар на таблицата: зависи ли времето за сходимост от диаметъра и по какъв
закон, и каква част от него се пада на откриването на отказа}

\subsection{Преходни цикли}

Броят и продължителността на преходните цикли за същите комбинации са дадени
\tabref{tab:loops}.

\begin{table}[H]
\caption{Преходни цикли при пренастройване на маршрутите}
\label{tab:loops}
\centering
\fontsize{11}{13}\selectfont
\begin{tabular}{lccc}
\toprule
Топология & Брой цикли & Средна продължителност & Максимална продължителност \\
\midrule
\TODO{} & \TODO{} & \TODO{} & \TODO{} \\
\TODO{} & \TODO{} & \TODO{} & \TODO{} \\
\bottomrule
\end{tabular}
\end{table}

\TODO{коментар на таблицата: при кои топологии изобщо възникват цикли и дали
продължителността им следва времето за сходимост}

\subsection{Сверяване с еталонната реализация}

\TODO{резултат от проверката за съвместимост с еталонната реализация: установява ли се
съседство, съвпадат ли изчислените маршрути, и къде има разминаване}

\subsection{Обсъждане}

\TODO{обсъждане на резултатите след попълване на таблиците: кои очаквания се потвърждават,
кои не, и какви са ограниченията на измерването}
